\chapter{Resultados}

El Sprint 1 entrega el primer resultado demostrable: desde un navegador se puede introducir un usuario y contraseña, autenticar contra la API real y acceder a una zona que no existe para visitantes anónimos. El rol aparece de forma comprensible y altera la navegación. La creación del primer administrador queda separada del uso diario y se bloquea una vez cumplida su función.

El incremento también reduce riesgo futuro. El modelo de datos completo hace visibles las relaciones y restricciones antes de construir catálogos; el contrato de error evita tratamientos distintos por módulo; la migración, los contenedores y la CI establecen un camino repetible que los sprints posteriores ampliarán.

Los Sprints 2 a 4 convierten ese esqueleto en una herramienta operativa demostrable. Ya se mantienen vehículos, conductores y clientes sin perder los registros dados de baja, y se crean, consultan y editan envíos asociados a esos datos. El listado paginado y sus filtros persistentes permiten localizar trabajo por estado, fecha y recursos; la normalización UTC evita que la misma operación cambie de significado entre navegador, API y PostgreSQL.

El Sprint 4 aporta por primera vez la lógica operacional central: los recursos se asignan sin reutilizar los que ya están ocupados, una capacidad insuficiente se comunica sin impedir el trabajo y el envío avanza por un ciclo irreversible con fechas reales. El Flujo 3 ya puede ejecutarse desde el alta hasta entrega o cancelación, salvo el registro de eventos que corresponde al siguiente incremento. La evidencia sobre PostgreSQL real confirma que las reglas no dependen del proveedor en memoria usado por las pruebas rápidas.

Todavía no procede evaluar la utilidad operacional completa: faltan los eventos, la administración y los indicadores. Esa evaluación se actualizará con métricas de cobertura funcional, defectos y resultados de los cuatro flujos al final del Sprint 7.
